% from ega2-1-fr.tex
\oldpage[II]{185}
et par définition $\pi_*(\mathcal{O}_V/\mathcal{J}^{n+1})$ est la fibre de
$\mathcal{O}_V/\mathcal{J}^{n+1}$ au point $j(\xi)$, c'est-à-dire, en posant
$C=\mathcal{O}_{j(\xi)}$ et en désignant par $\mathfrak{m}$ l'idéal maximal
de $C$, le $C$-module $C/\mathfrak{m}^{n+1}$~; l'assertion (ii) est alors
évidente.

Pour démontrer (i), notons que la question est locale sur $X$~; on peut donc
se restreindre au cas où $X$ est affine. Soit $U$ un ouvert affine dans $V$~;
$j(X)\cap U$ est un ouvert affine dans $j(X)$, donc
$U_0=\pi(j(X)\cap U)$, qui lui est isomorphe, un ouvert affine dans $X$~;
pour tout ouvert affine $W_0\subset U_0$ dans $X$,
$W=\pi^{-1}(W_0)\cap U$ est un ouvert affine dans $V$, puisque $X$ est un
schéma (\hyperref[I.5.5.10-fr]{I, 5.5.10})~; en particulier
$U'=U\cap\pi^{-1}(U_0)$ est un ouvert affine dans $V$ et on a évidemment
$\pi(U')=U_0$ et $j(U_0)=j(X)\cap U$. Cela étant, par définition
\[
  \Gamma(W_0,\pi_*(\mathcal{O}_V/\mathcal{J}^{n+1}))
  =\Gamma(\pi^{-1}(W_0),\mathcal{O}_V/\mathcal{J}^{n+1})~;
\]
mais comme tout point de $\pi^{-1}(W_0)$ n'appartenant pas à $j(W_0)$ a un
voisinage ouvert dans $\pi^{-1}(W_0)$ ne rencontrant pas $j(X)$ et dans
lequel $\mathcal{O}_V/\mathcal{J}^{n+1}$ est donc nul, il est clair que les
sections de $\mathcal{O}_V/\mathcal{J}^{n+1}$ au-dessus de $\pi^{-1}(W_0)$ et
au-dessus de $W$ se correspondent biunivoquement. Autrement dit, si $\pi'$ est
la restriction de $\pi$ à $U'$, les $(\mathcal{O}_X|U_0)$-Modules
\[
  \pi_*(\mathcal{O}_V/\mathcal{J}^{n+1})|U_0
  \quad\hbox{et}\quad
  \pi'_*\bigl((\mathcal{O}_V/\mathcal{J}^{n+1})|U'\bigr)
\]
sont identiques. Comme $U'$ et $U_0$ sont affines et que les $U_0$ recouvrent
$X$, on en conclut (\hyperref[I.1.6.3-fr]{I, 1.6.3}) que
$\pi_*(\mathcal{O}_V/\mathcal{J}^{n+1})$ est quasi-cohérent, et la
démonstration est identique pour
$\pi_*(\mathcal{J}/\mathcal{J}^{n+1})$.

Pour démontrer enfin (iii), notons que $\mathcal{S}$ n'est autre que
$\mathbf{S}_{\mathcal{O}_X}(\mathcal{L})$~; on déduit donc canoniquement de
$\varphi$ un homomorphisme de $\mathcal{O}_X$-Algèbres
\[
  \psi:\mathcal{S}\longrightarrow
  \varprojlim_n\pi_*(\mathcal{O}_V/\mathcal{J}^{n+1})
\]
(\hyperref[II.1.7.4-fr]{1.7.4})~; en outre, cet homomorphisme applique
$\mathcal{L}^{\otimes n}$ dans
$\varprojlim_m\pi_*(\mathcal{J}^n/\mathcal{J}^{m+1})$, donc est continu pour
les topologies considérées, et se prolonge bien par suite en un homomorphisme
\[
  \widehat{\psi}:\widehat{\mathcal{S}}\longrightarrow
  \varprojlim_n\pi_*(\mathcal{O}_V/\mathcal{J}^{n+1}).
\]
Pour voir qu'il s'agit d'un isomorphisme, on peut, comme dans la démonstration
de (i), se réduire au cas où $X=\operatorname{Spec}(A)$ et
$V=\operatorname{Spec}(B)$ sont affines, avec
$\mathcal{J}=\widetilde{\mathfrak{J}}$, où $\mathfrak{J}$ est un idéal de
$B$~; il correspond à $\pi$ une injection $A\to B$, identifiant $A$ à un
sous-anneau de $B$ supplémentaire de $\mathfrak{J}$, et $\mathcal{L}$ (resp.
$\pi_*(\mathcal{O}_V/\mathcal{J}^{n+1})$) est le $\mathcal{O}_X$-Module
quasi-cohérent associé au $A$-module $L=\mathfrak{J}/\mathfrak{J}^2$ (resp.
$B/\mathfrak{J}^{n+1}$). Comme $\mathcal{J}$ est un
$\mathcal{O}_V$-Module \emph{inversible}, on peut en outre supposer que
$\mathfrak{J}=Bt$, où $t$ est non diviseur de 0 dans $B$. De la relation
$B=A\oplus Bt$ on déduit alors, pour tout $n>0$,
\[
  B=A\oplus At\oplus At^2\oplus\cdots\oplus At^n\oplus Bt^{n+1}
\]
donc on a un $A$-isomorphisme canonique de l'anneau de séries formelles
$A[[T]]$ sur $C=\varprojlim B/\mathfrak{J}^{n+1}$ faisant correspondre $t$ à
$T$. D'autre part, on a $L=A\bar t$, où $\bar t$ est la classe de $t$
mod. $Bt^2$, et l'homomorphisme $\varphi$ applique par hypothèse $\bar t$ sur
un élément $t'\in C$ congru à $t$ mod. $Ct^2$. On en déduit par récurrence
sur $n$ que
\[
  A\oplus At'\oplus\cdots\oplus At'{}^n\oplus Ct^{n+1}
  =A\oplus At\oplus\cdots\oplus At^n\oplus Ct^{n+1}
\]
ce qui prouve que l'homomorphisme $\widehat{\psi}$ correspond bien à un
isomorphisme de $\prod_{n\geqslant0}L^{\otimes n}$ sur $C$.
\end{proof}

\begin{lemma}[8.10.2]
\label{II.8.10.2-fr}
Sous les hypothèses de (\hyperref[II.8.10.1-fr]{8.10.1}), soit
$g:X'\to X$ un morphisme,

